% !TEX root = ../main.tex
\section{Costly Contracting}\label{thr:costly-contracting}

\paragraph{Definisjon.} Costly contracting (\emph{kostbar kontrahering}) er forutsetningen om at det å skrive, forhandle, observere og håndheve kontrakter er ressurskrevende, og at kontrakter derfor i praksis er \emph{ufullstendige}: de kan ikke spesifisere hver mulig fremtidig tilstand. Konsekvensen er at insentivkonflikter ikke kan elimineres ved kontraktdesign alene; det optimale er en \emph{second-best}-løsning der residuelle tap aksepteres mot at kontraktomkostningene holdes nede. [SLIDE]

\subsection*{Forklaring}

I en idealisert verden av \thref{costless-contracting}{costless contracting} kan partene gratis skrive en \emph{komplett kontrakt} som binder agenten til å handle i prinsipalens interesse i alle tenkelige tilstander. Da forsvinner alle de fem incitamentkonfliktene mellom eier og leder via kontraktdesign og førsteorden-løsningen er \emph{first-best}. Denne benchmark-en brukes i Q03. [SLIDE/BOK]

I virkeligheten påløper det \emph{transaction costs} (Coase 1937, Williamson 1979) ved hver fase av kontraheringen:
\begin{itemize}\setlength{\itemsep}{2pt}
\item \textbf{Search and information costs} -- finne motpart og innhente informasjon om type/kvalitet (her treffer \thref{adverse-selection}{adverse selection}).
\item \textbf{Bargaining and decision costs} -- forhandle kontraktsvilkår, prise inn risiko og asymmetri.
\item \textbf{Drafting costs} -- formulere kontrakten skriftlig, inkludert tilstandsavhengighet.
\item \textbf{Policing/monitoring costs} -- observere agentens handling i ettertid (her treffer \thref{moral-hazard}{moral hazard}).
\item \textbf{Enforcement costs} -- bruke domstoler eller voldgift for å håndheve brudd.
\end{itemize}

I tillegg kommer \emph{begrenset rasjonalitet} (Simon): partene kan ikke forutse alle fremtidige tilstander. Dermed blir kontraktene nødvendigvis \emph{ufullstendige}. Dette åpner for \emph{renegotiation}: når en uforutsett tilstand inntreffer, må partene forhandle på nytt -- og den parten som har størst makt på det tidspunktet kan utøve \emph{hold-up}. [BOK]

Costly contracting er derfor selve forutsetningen som gjør de fem incitamentkonfliktene reelt problematiske: under costless contracting kunne kontrakten skrevet dem bort. Under costly contracting velger man en \emph{second-best}-løsning og aksepterer \emph{agency costs} (\thref{agency-costs}{monitoring + bonding + residual loss}). Den optimale insentivkontrakten balanserer:
\begin{itemize}\setlength{\itemsep}{2pt}
\item Marginal verdi av økt motivasjon (mer $\beta$, lavere shirking) mot
\item Marginal kost av økt risikopremie (agenten bærer mer av $\mu$-støy) og økt monitoring/bonding.
\end{itemize}

To-virksomhets-perspektivet i Q04 (prinsipalen er én virksomhet, agenten en annen virksomhet) flytter problemet fra ansettelseskontrakt til \emph{interfirm contracting}: utviklingsavtaler, leverandørkontrakter, joint ventures, lisensavtaler. Her er adverse selection særlig tydelig (du kjenner ikke partnerens reelle kompetanse) og moral hazard tar form av f.eks.\ underleveranse, fortielse av problemer, fri-ridning på fellesinvesteringer. Williamson-rammeverket -- spesifikke aktiva, hold-up, vertikal integrasjon -- hører hjemme her og brukes utfyllende i Q11/Q12. [BOK]

\subsection*{Kjerneforutsetninger}
\begin{itemize}
\item \textbf{Begrenset rasjonalitet:} Partene kan ikke forutse alle fremtidige tilstander.
\item \textbf{Asymmetrisk informasjon:} Pre- og post-kontraktuelt (adverse selection og moral hazard).
\item \textbf{Opportunisme:} Aktørene utnytter informasjonsfortrinn og hull i kontrakten når det lønner seg.
\item \textbf{Verifiserbarhetsproblem:} Tredjepart (domstol) kan ikke verifisere alle relevante variabler -- noen ting er observerbare for partene, men ikke verifiserbare. Disse kan ikke kontraktreguleres.
\item \textbf{Spesifikke investeringer skaper hold-up-risiko:} Når en part investerer i aktiva med liten verdi utenfor kontrakten, kan motparten reforhandle vilkårene ex post.
\end{itemize}

\subsection*{Muligheter (hva forklarer/løser teorien?)}
\begin{itemize}
\item Forklarer hvorfor virksomheter eksisterer (Coase): noen transaksjoner er for dyre å koordinere via kontrakter i markedet -- de internaliseres i firmaet.
\item Forklarer behovet for monitoring og bonding -- og hvorfor disse koster.
\item Begrunner valget mellom \emph{make or buy} (vertikal integrasjon vs.\ outsourcing) basert på relative kontraktomkostninger.
\item Forklarer hvorfor man bruker langsiktige relasjoner, reputasjon og implisitte (ikke-rettslig håndhevbare) kontrakter.
\item Gir grunnlag for residual control rights -- den som eier aktiva har myndighet i tilstander kontrakten ikke har spesifisert (Grossman-Hart-Moore).
\item Begrunner second-best-tankegangen: man kan ikke løse alle insentivproblemer, bare optimalisere balansen mellom motivasjon, risikodeling og kontraktomkostninger.
\end{itemize}

\subsection*{Begrensninger og kritikk}
\begin{itemize}
\item Vanskelig å måle transaksjonsomkostninger empirisk -- mange er implisitte.
\item Modellen sier sjelden \emph{hvor mye} kontraktomkostninger må være før vertikal integrasjon lønner seg.
\item Forutsetningen om opportunisme kan overdrives -- relasjonelle og tillitsbaserte kontrakter fungerer bedre enn modellen forutsier (skandinavisk arbeidsliv, japanske keiretsu).
\item Antar at intern organisering (firmaet) er billigere enn marked når kontraktomkostningene blir høye -- men også interne hierarkier har bureaucracy costs (jf.\ Williamson selv).
\item Skiller ikke alltid skarpt mellom kontraktomkostninger og rene insentivkostnader -- agency costs er en del av kontraktomkostningene.
\end{itemize}

\subsection*{Modell / illustrasjon}

\textbf{First-best vs.\ second-best under costly contracting:}

\begin{center}
\begin{tabular}{@{}p{4.0cm}p{5.0cm}p{5.5cm}@{}}
\toprule
 & \textbf{Costless contracting (Q03)} & \textbf{Costly contracting (Q04)} \\
\midrule
Informasjon & Symmetrisk eller gratis verifiserbar & Asymmetrisk (adverse selection + moral hazard) \\
Kontrakt & Fullstendig, statsavhengig & Ufullstendig \\
Insentivkonflikter & Skrives bort & Reduseres, ikke elimineres \\
Optimal $\beta$ & Setter førstebeste innsats $e^*$ & Trade-off motivasjon vs.\ risikopremie \\
Resultat & First-best & Second-best $+$ agency costs \\
Verktøy & Kontraktdesign alene & Screening, signaling, monitoring, bonding, reputasjon, integrasjon \\
\bottomrule
\end{tabular}
\end{center}

Forelesnings\-eksemplet (BSZ kap. 10): med contracting costs 800, monitoring 400, bonding 400, residual loss 100, blir transaksjonens nettoverdi $S' = 2{,}500 - 800 - 100 = 1{,}600$ -- mot $2{,}500$ under costless contracting. Differansen er det totale agency-tapet ved costly contracting. [SLIDE]

\subsection*{Eksempler}
\begin{itemize}
\item \textbf{Equinor og lisensavtaler på sokkelen:} Detaljerte joint operating agreements (JOA) mellom partnere på en lisens er enormt omfattende (hundrevis av sider) nettopp fordi kontraktomkostningene knyttet til langsiktig samarbeid med flere parter er høye, og fordi kontraktene må håndtere uforutsette geologiske og kommersielle tilstander. [EKS]
\item \textbf{Statoil-Hydro-fusjonen (2007):} Måneder med due-diligence er search and information cost; integrasjonsprosessen etterpå er forhandlings- og policing-kostnad. Selv etter fusjonen var det betydelige residuelle insentivproblemer mellom kulturer og divisjoner. [EKS]
\item \textbf{Mærsk og leverandørkontrakter for elektrolyseanlegg:} Kontrakter med teknologileverandører for Power-to-X må spesifisere ytelse, garantier og milepæler -- klassisk costly contracting med stor adverse-selection-risiko (er teknologien moden?) og moral hazard (vil leverandøren skyve problemer foran seg?). [CASE]
\item \textbf{IKEA-leverandører i Polen og Vietnam:} IKEA bruker langsiktige rammeavtaler med detaljerte kvalitets- og leveringskrav, plus stedlige inspektører (monitoring) -- second-best-løsning for å holde leverandørrelasjoner verdiskapende uten vertikal integrasjon. [EKS]
\end{itemize}

\subsection*{Kobling til søyle og andre teorier}
Søyle: \SThree\ primært (insentivkontrakter under costly contracting), men også \SOne\ (residuelle beslutningsrettigheter ved kontraktufullstendighet) og \STwo\ (monitoring som prestasjonsmåling).

Relaterte teorier: \thref{costless-contracting}{Costless Contracting} (motsatsen, brukt i Q03), \thref{principal-agent}{Principal-Agent}, \thref{moral-hazard}{Moral Hazard}, \thref{adverse-selection}{Adverse Selection}, \thref{agency-costs}{Agency Costs}, \thref{incentive-conflicts-owner-manager}{Incitamentkonflikter}, \thref{specific-vs-general-knowledge}{Spesifik vs.\ generell viden}.

\paragraph{Brukt i spørsmål:} Q04 (primært), berøres i Q11/Q12 (vertikal integrasjon vs.\ outsourcing).
